% ============================================================
%  Conclusion
% ============================================================

\section{Conclusion}
\label{sec:conclusion}

This document has presented the algorithms, data structures, and
engineering decisions underlying \FractalShark{}, a GPU-accelerated
Mandelbrot set renderer capable of interactive exploration at zoom
depths exceeding $10^{10000}$.  \Cref{fig:system-architecture}
provides a high-level view of the pipeline.

\begin{figure}[!htbp]
\centering
\begin{tikzpicture}[
  every node/.style={font=\small},
  block/.style={draw, rounded corners, minimum height=0.65cm,
                minimum width=2.2cm, align=center},
  comp/.style={block, fill=blue!10},
  io/.style={block, fill=orange!10},
  infra/.style={block, fill=gray!12, dashed},
  arr/.style={-{Stealth[length=5pt]}, thick},
  label/.style={font=\scriptsize\bfseries, anchor=north west},
  grp/.style={draw, dashed, rounded corners, inner sep=6pt, gray!70}
]

% --- Row 1: Input ---
\node[io] (params) {View Parameters};
\node[comp, right=0.6cm of params] (pzbb) {PointZoom-\\BBConverter};
\node[comp, right=0.6cm of pzbb] (deltas) {Per-pixel\\deltas};

\draw[arr] (params) -- (pzbb);
\draw[arr] (pzbb) -- (deltas);

% --- Row 2: Core computation ---
\node[comp, below=1.0cm of params] (reforbit) {Reference Orbit\\(RefOrbitCalc)};
\node[comp, right=0.6cm of reforbit] (la) {LA Setup\\(LAReference)};
\node[comp, right=0.6cm of la] (perturb) {Perturbation\\Kernels};

\draw[arr] (deltas) -- ++(0,-0.45) -| (perturb.north);
\draw[arr] (reforbit) -- (la);
\draw[arr] (la) -- (perturb);

% --- Row 3: Output ---
\node[comp, below=1.0cm of reforbit] (aa) {Antialiasing};
\node[comp, right=0.6cm of aa] (color) {Palette\\Coloring};
\node[io, right=0.6cm of color] (gl) {GL Display};

\draw[arr] (perturb) -- ++(0,-0.45) -| (aa.north);
\draw[arr] (aa) -- (color);
\draw[arr] (color) -- (gl);

% --- Side: GPU arithmetic ---
\node[infra, left=0.6cm of reforbit] (gpu) {GPU Arith.\\(HpSharkFloat)};
\draw[arr] (gpu) -- (reforbit);

% --- Side: Feature Finder ---
\node[infra, below=0.35cm of gpu] (ff) {Feature\\Finder};
\draw[arr, dashed] (ff.east) -- ++(0.3,0) |- ([yshift=-2pt]reforbit.south west);

% --- Side: Memory ---
\node[infra, below=0.35cm of ff] (mem) {Growable-\\Vector};
\draw[arr, dashed] (mem.east) -- ++(0.6,0) |- ([yshift=2pt]reforbit.south);

% --- Group labels ---
\begin{scope}[on background layer]
  \node[grp, fit=(pzbb)(deltas)(perturb), label={[label]above left:Rendering}] {};
  \node[grp, fit=(reforbit)(la), label={[label]above left:Ref.\ Orbit}] {};
  \node[grp, fit=(gpu)(ff)(mem), label={[label]above left:Infrastructure}] {};
\end{scope}

\end{tikzpicture}
\caption[End-to-end rendering pipeline]{High-level architecture of the
  \FractalShark{} rendering pipeline.  Solid arrows denote the primary
  data flow; dashed arrows indicate supporting subsystems.  Grouped
  regions correspond to the document's Part structure.}
\label{fig:system-architecture}
\end{figure}

The exposition began with the direct rendering kernelsthat map pixels
to the complex parameter plane and evaluate the escape-time iteration
(\cref{sec:mapping}), followed by the numeric type hierarchy---from
native \code{float} through multi-limb \code{HpSharkFloat}---that
provides the precision each stage requires (\cref{sec:hdr-float}).
Building on this foundation, perturbation theory
(\cref{sec:perturbation-concept}) reduces per-pixel work from
arbitrary-precision arithmetic to lightweight floating-point deltas,
and linear approximation (\cref{sec:approx-accel}) eliminates large
blocks of iterations entirely by replacing them with a single
matrix--vector product.

The second half of the document addressed the computation of the
reference orbit that perturbation demands
(\cref{sec:ref-orbit-calc}), the NTT-based GPU arithmetic engine that
accelerates it (\cref{subsec:ref-orbit-gpu}), and the Feature Finder
that applies the same high-precision machinery to automatic detection
of periodic points (\cref{sec:feature-finder}).  Supporting
infrastructure---post-iteration processing and colour mapping
(\cref{sec:misc-details}), together with the disk-backed growable
vectors that manage large reference orbits
(\cref{sec:disk_backed_growable_vectors})---completes the rendering
pipeline.

Several design decisions merit emphasis.  \FractalShark{} adopts
\emph{rebasing} rather than glitch correction to handle loss of
significance during perturbation iteration, avoiding the need for
post-hoc detection and re-rendering of glitched pixels.  Linear
approximation version~2 (LA~v2) was chosen over series approximation
for its period-aligned skip structure, which enables longer effective
jumps at deep zooms.  NTT-based multiplication was preferred over
Karatsuba for the GPU arithmetic backend, as the former maps naturally
to the highly parallel CUDA execution model.  Throughout the
codebase, C++ template instantiation over a fixed parameter-set
lattice allows the compiler to specialise kernels for each combination
of precision, sub-type, and algorithmic mode while keeping the source
generic.

Taken together, these techniques form a layered architecture in which
each subsystem can be understood, tested, and tuned independently.
\cref{sec:future-directions} surveyed several alternative and
complementary techniques---including series approximation, BLA,
super-resolution, and histogram-based colouring---that represent
natural extensions of the current design.

\FractalShark{} remains under active development.  As new algorithmic
ideas emerge from the deep-zoom rendering community and as GPU
hardware continues to evolve, the architecture described here provides
a foundation for continued experimentation and improvement.
